\noindent\textbf{Stage-2 selection and the scope of this appendix.}
No class-C or class-D item justified a bounded research development in this run.
The following elementary proofs support the critical edition of already identified
material. They verify source readings and reconstructions, with no historical
novelty claim. The item records retain development status \textit{annotated};
none is reclassified as an open problem to justify this appendix.


% Critical checks of source formulas, not Stage 2 research developments.
\subsection{The sign distribution and its moments}
The conventions change between P000503 and P000505. We use the latter: for an odd
prime \(p\), put \(N=p^m\) and, for \(0\le a<N\), set
\[
 \mu(a+N\mathbb Z_p)=\begin{cases}-1,&2a<N,\\+1,&2a>N.\end{cases}
\]
The following verifies the notebook formula and its normalization; it is not a
claim to a new theorem on \(p\)-adic integration. The underlying regularization is
part of \citet[\S8.4]{manin1973periods}.

\begin{proposition}[Moment correction for items 1973-303.14 and 1973-304.2--3]\ours\quad
These values define a bounded clopen distribution. With Bernoulli numbers
\(B_j\) defined by \(t/(e^t-1)=\sum_{j\ge0}B_jt^j/j!\), its polynomial moments are
\[
 I_k:=\int_{\mathbb Z_p}x^k\,d\mu
   =\frac{(4-2^{1-k})B_{k+1}}{k+1},\qquad k\ge0.
\]
Moreover \(\int_{\mathbb Z_p^\times}x^k\,d\mu=(1-p^k)I_k\).
\end{proposition}
\begin{proof}
Write \(p=2r+1\). Among the children \(a+jN\), \(0\le j<p\), the first \(r\)
lie below \(pN/2\), the last \(r\) lie above it, and the middle child lies below
exactly when \(a<N/2\). The first and last children cancel in the sum; the middle
child has the parent's sign. Thus the distribution compatibility equation holds
at every level, including level zero. Every residue-ball value has absolute value one, so every clopen value has absolute value at most one by the non-Archimedean triangle inequality.
Uniform approximation by locally constant functions therefore defines its
integral on continuous \(\mathbb Q_p\)-valued functions.

For the polynomial \(x^k\), the canonical residue-tag sum is
\[
 S_k(N)=\sum_{a=0}^{N-1}a^k-2\sum_{a=0}^{(N-1)/2}a^k,
\]
with the convention \(a^0=1\), also at \(a=0\). Let
\(B_j(X)\) be the Bernoulli polynomial. Its elementary difference identity
\(B_{k+1}(X+1)-B_{k+1}(X)=(k+1)X^k\) gives the exact rational identity
\[
 S_k(N)=\frac{B_{k+1}(N)-2B_{k+1}((N+1)/2)+B_{k+1}(0)}{k+1}.
\]
As \(m\to\infty\), \(N=p^m\to0\) in \(\mathbb Q_p\). Polynomial continuity yields
\[
 I_k=\frac{2(B_{k+1}-B_{k+1}(1/2))}{k+1}.
\]
The Bernoulli duplication identity gives
\(B_{k+1}(1/2)=(2^{-k}-1)B_{k+1}\), which proves the displayed result.
In particular \(I_0=-1\), \(I_1=1/4\), \(I_3=-1/32\), and \(I_{2r}=0\) for
\(r\ge1\). These values alone already distinguish the source's convention from
several possible Bernoulli-index corrections.

Finally \(\mu(pU)=\mu(U)\) for each residue ball \(U\): the defining inequality
is unchanged when both its numerator and modulus are multiplied by \(p\).
Taking finite sums and then limits gives
\(\int_{p\mathbb Z_p}x^k\,d\mu=p^kI_k\); subtracting from the total proves the
unit formula. The opposite convention on P000503 negates every moment.
\end{proof}

\paragraph{Independent verification of the moment calculation.}
\oursComputed\ The accompanying \texttt{code/independent\_checks.py} derives
Bernoulli numbers recursively using exact rational arithmetic. It checks the
finite-sum identity for \(p=3,5,7\), \(k=0,\ldots,9\), and \(m=1,\ldots,4\):
all 120 comparisons agree exactly. The recorded \(p\)-adic error valuations
are descriptive checks; the proof of the limit is polynomial continuity,
not extrapolation from those cases. The sign convention, duplication
identity and unit restriction were independently rederived, with the same
formula and the values \(-1,1/4,-1/32\) at \(k=0,1,3\).


\subsection{Dependence on tags for the invariant distribution}
\begin{proposition}[Item 1973-303.16]\ours\quad
The translation-invariant distribution of total mass one on \(\mathbb Z_p\)
does not have a tag-independent Riemann integral of \(x\).
\end{proposition}
\begin{proof}
On level \(m\) it gives each residue ball mass \(p^{-m}\). With canonical tags,
\[
 p^{-m}\sum_{a=0}^{p^m-1}a=(p^m-1)/2\longrightarrow-1/2.
\]
Change the tag in just the ball \(p^m\mathbb Z_p\) from \(0\) to \(p^m\). This
adds \(p^mp^{-m}=1\) at every level and changes the limit to \(1/2\). Both
partitions have mesh tending to zero; only their tags differ. This does not
contradict integration with a separately fixed representative convention.
\end{proof}

\subsection{Two checks on categorical realization}
\begin{example}[A nondegenerate simplex with a degenerate face, item 1973-308.5]\ours\quad
In the nerve of the group \(C_2=\{1,g\}\), the two-simplex \((g,g)\) is
nondegenerate, since degeneracies insert an identity arrow. Its middle face is
\(g^2=1\), the degenerate one-simplex at the unique vertex. Thus the disputed
sentence on P000515 is false even for the elementary group example on P000516.
The usual realization quotient still attaches a cell for every nondegenerate
simplex; its attaching map may identify part of its boundary with lower cells.
\end{example}

\begin{proposition}[Two Q-construction checks, items 1973-311.11--12]\ours\quad
Let \(A\) be a unital ring and \(P(A)\) its exact category of finitely generated
projective modules. The map \(BGL_n(A)\to BQP(A)\) obtained by viewing
\(A^n\) and its automorphisms in \(QP(A)\) is nullhomotopic. Direct sum with
\(A\) induces a map \(BQP(A)\to BQP(A)\) homotopic to the identity.
\end{proposition}
\begin{proof}
For the first assertion, a Q-morphism \(0\to A^n\) is represented by
\(0\xleftarrow{\mathrm{id}}0\rightarrowtail A^n\). Postcomposing it with any
automorphism of \(A^n\) gives the same Q-morphism. Hence these arrows form a
natural transformation from the constant zero functor on the one-object
category \(GL_n(A)\) to the automorphism functor. Natural transformations realize
to homotopies, by \citet[\S1, Proposition~2]{quillen1973higher}.

For the second assertion, let \(S(X)=X\oplus A\), and define S on a span by
summing the identity of A on both arrows. The admissible inclusion
\(X\hookrightarrow X\oplus A\) gives a Q-morphism represented by
\(X\xleftarrow{\mathrm{id}}X\rightarrowtail X\oplus A\). To verify naturality
for \(X\xleftarrow{p}E\xrightarrow{i}Y\), pull back
\(p\oplus1:E\oplus A\to X\oplus A\) along the inclusion of X. This pullback
is E, so either composite is represented by
\(X\xleftarrow{p}E\xrightarrow{(i,0)}Y\oplus A\). Thus these components give
\(\mathrm{id}\Rightarrow S\), proving the required homotopy.

The first natural transformation only concerns the automorphism subcategory.
It does not extend to a contraction of Q: pulling a general admissible
epimorphism back along zero can have a nonzero kernel. Neither assertion replaces
the substantive comparison theorem
\(\Omega BQP(A)\simeq K_0(A)\times BGL(A)^+\), checked in
\citet[Corollary~IV.7.2]{weibel2013kbook}.
\end{proof}

\subsection{Critical verification of the conic example (1973-322.9)}
\label{sec:conic-verification-1973}

This is a verification appendix to the settled complex formulation indexed as
GP-1973-32. It is not a Stage~2 contribution to a C or D problem. The explicit
certificates below are proved here to resolve a consequential ambiguity in the
notes; no historical priority or novelty is claimed.

\paragraph{What the source supports.}
P000535 first gives an example explicitly over \(\mathbb R\), with the
unimodular triple \((x,y,1-x^2-y^2)\) and the comment that degree obstructs its
reduction. After an undefined winding symbol and the definition of symplectic
\(K_1\), the page asks about \(z_1^2+z_2^2=1\), entire functions, and entire
functions of exponential growth. The two resulting functions are required to
have \ru{не имеют общ. делит.}, literally no common divisor. The parenthesis
\ru{в полиномах нет, в непр. ф-ях есть} says that polynomials do not suffice
but continuous functions do. These words and the equation are legible in the
scan; the coefficient field is not explicitly restated.

The strongest coherent existence question furnished by the displayed formulas
is the following, once the algebra is specified. Put
\[
 c=x^2+y^2-1,
 \qquad g_1=x+cf_1,\qquad g_2=y+cf_2.
\]
Can \(f_1,f_2\) be chosen so that \(g_1,g_2\) generate the unit ideal in the
chosen function algebra? This is stronger than the literal common-divisor
condition. In \(\mathbb C[x,y]\), taking \(f_1=f_2=0\) already leaves the
relatively prime pair \(x,y\), although that pair has a common zero. The
stable-rank context is the reason for reconstructing unimodularity. It is not
an independent reading of an unrecorded condition. The initial triple really
is unimodular, since \(x\cdot x+y\cdot y-c=1\).

\paragraph{The verified literature and its scope.}
The real obstruction is an old one: Vaserstein's 1971 Theorem~8 and its proof
use the same sphere triple to establish the full polynomial stable rank over
subfields of \(\mathbb R\) \citep[Russian p.~27]{vaserstein1971stable}.
For unrestricted holomorphic functions, Brudnyi's theorem gives
\[
 \operatorname{sr}\mathcal O(X)
   =\left\lfloor\frac{\dim_{\mathbb C}X}{2}\right\rfloor+1
\]
for finite-dimensional Stein spaces, and therefore
\(\operatorname{sr}\mathcal O(\mathbb C^2)=2\)
\citep[Theorem~1.1]{brudnyi2019stein}. Its preprint circulated in March 2018;
the journal article appeared in January 2019. This theorem applies to every
unimodular triple of entire functions on \(\mathbb C^2\), but supplies no
exponential-growth restriction. The following particular example needs neither
that general theorem nor a new open-problem designation.

\begin{proposition}[A polynomial certificate over the complex numbers]\ours\quad
Let \(c=x^2+y^2-1\) and set
\[
 f_1=\frac12,\qquad
 f_2=\frac{i(x+2)-y}{2},\qquad
 a=x+cf_1,\qquad b=y+cf_2.
\]
Then \(a,b\) generate the unit ideal in \(\mathbb C[x,y]\). More precisely,
with
\[
 s=x+1-iy,\qquad
 A=s(a+1)-a-2,\qquad B=\frac{i s^2}{2},
\]
one has \(Aa+Bb=1\) identically.
\end{proposition}

\begin{proof}
Put \(t=x+1+iy\). Direct expansion gives
\[
 ts=(x+1)^2+y^2=2(a+1),
 \qquad
 b=\frac{t^2}{2i}+ita.
\]
For completeness, the second identity follows by expanding its right side:
its coefficient of \(i\), after grouping over \(\mathbb Q[x,y]\), is
\(i((x+1)a-((x+1)^2-y^2)/2)\), and its remaining part is
\((x+1)y-ya\); substituting \(a=x+c/2\) gives
\(y+c(i(x+2)-y)/2\). Hence
\[
 Bb=\frac{s^2t^2}{4}-\frac{s^2t a}{2}
    =(a+1)^2-s(a+1)a.
\]
Combining this with \(Aa=s(a+1)a-a^2-2a\) proves
\[
 Aa+Bb=(a+1)^2-a^2-2a=1.
\]
Every coefficient belongs to \(\mathbb Q(i)\), so this is an exact polynomial
identity. In particular \(a,b\) have no simultaneous zero anywhere in
\(\mathbb C^2\).
\end{proof}

\begin{corollary}[Entire and exponential-type formulations]\ours\quad
The complex conic triple is reducible in the polynomial algebra, in
\(\mathcal O(\mathbb C^2)\), and in the algebra
\[
 \mathcal E(\mathbb C^2)=
 \{f\in\mathcal O(\mathbb C^2):
   \exists C,T>0\ \forall z\in\mathbb C^2,   |f(z)|\le C e^{T(|z_1|+|z_2|)}\}.
\]
The reduction and its Bezout coefficients can all be taken to be polynomials.
\end{corollary}

\begin{proof}
The certificate persists under either inclusion of polynomial algebras.
Every polynomial \(p\) satisfies
\(|p(z)|\le C_\epsilon e^{\epsilon(|z_1|+|z_2|)}\) for every
\(\epsilon>0\), because a fixed power of \(1+r\) times
\(e^{-\epsilon r}\) is bounded on \([0,\infty)\). The algebra
\(\mathcal E\) is closed under addition and multiplication, with the
constants and types allowed to depend on the function. Thus both generators
and the explicit Bezout coefficients lie in it.
\end{proof}

\begin{proposition}[An entire certificate keeping the first coordinate]\ours\quad
There is also a solution with \(f_1=0\). Define
\[
 h(w)=\exp\!\left(\frac{i\pi(w-1)}2\right),\qquad
 F(w)=\frac{h(w)-w}{w^2-1},
\]
where the two removable values are filled in. Then
\[
 f_2(x,y)=F(y),\qquad
 g_1=x,\qquad g_2=h(y)+x^2F(y),
\]
and
\[
 -\frac{xF(y)}{h(y)}\,g_1+\frac1{h(y)}\,g_2=1.
\]
All the displayed entire functions belong to \(\mathcal E(\mathbb C^2)\).
\end{proposition}

\begin{proof}
The values \(h(1)=1\) and \(h(-1)=-1\) show that the numerator of \(F\)
vanishes at both roots of the denominator. Both roots are simple, and the
removable values obtained by differentiation are
\[
 F(1)=\frac{i\pi}{4}-\frac12,\qquad
 F(-1)=\frac{i\pi}{4}+\frac12.
\]
Consequently \(F\) is entire and
\[
 cF(y)=x^2F(y)+h(y)-y.
\]
This is precisely the required form for \(g_2=y+cf_2\). The function \(h\)
never vanishes. Substitution proves the Bezout identity for all \((x,y)\),
including \(y=\pm1\) by removability.

Here is a uniform growth check, not only a pointwise asymptotic argument.
Write \(\alpha=\pi/2\). For every \(w\),
\(|h(w)|,|h(w)^{-1}|\le e^{\alpha|w|}\). If \(|w|\ge2\), then
\[
 |w^2-1|\ge |w|^2-1\ge\tfrac34|w|^2,
 \qquad
 |F(w)|\le\frac{4(e^{\alpha|w|}+|w|)}{3|w|^2}.
\]
On \(|w|\le2\), the removable extension is bounded by compactness. Therefore
some finite \(K\) satisfies \(|F(w)|\le K e^{\alpha|w|}\) for all \(w\),
and \(|F(w)/h(w)|\le K e^{2\alpha|w|}\). Using
\(r\le e^r\) and \(r^2\le2e^r\) for \(r\ge0\) gives, for example,
\[
 \begin{split}
 |g_2(x,y)|&\le(1+2K)e^{|x|+\alpha|y|},\\
 |xF(y)/h(y)|&\le K e^{|x|+\pi|y|},\\
 |1/h(y)|&\le e^{\alpha|y|}.
 \end{split}
\]
Together with the bounds for \(g_1,f_1,f_2\), these imply a common estimate
\(C e^{(\pi+1)(|x|+|y|)}\) after increasing \(C\). No optimal type is
claimed.
\end{proof}

\begin{proposition}[The real-valued interpretation has the opposite answer]\ours\quad
Let \(D^2=\{(x,y)\in\mathbb R^2:x^2+y^2\le1\}\). For continuous
real-valued \(f_1,f_2\) on \(D^2\), the pair
\[
 (x+(x^2+y^2-1)f_1,\ y+(x^2+y^2-1)f_2)
\]
has a zero in \(D^2\). In particular there is no such reduction over
\(\mathbb R[x,y]\).
\end{proposition}

\begin{proof}
If this pair were nowhere zero, dividing it by its Euclidean norm would give
a continuous map \(r:D^2\to S^1\). On \(S^1\) the factor
\(x^2+y^2-1\) is zero, so \(r\) restricts to the identity. The composite
\(S^1\hookrightarrow D^2\xrightarrow{r}S^1\) would induce the identity on
\(H_1(S^1;\mathbb Z)=\mathbb Z\), but it factors through
\(H_1(D^2;\mathbb Z)=0\). This contradiction proves the assertion. Changing
\(c\) to \(-c\), as in the preceding notebook example, only changes the
signs of the coefficient functions.
\end{proof}

\paragraph{Independent checks.}
\oursComputed\ The accompanying \texttt{check\_conic.py} uses exact pairs of
rational numbers for coefficients in \(\mathbb Q(i)\), expands both
intermediate identities, and verifies that every coefficient of
\(Aa+Bb-1\) is zero. Its recorded output gives degrees
\(0,1,2,3,3,2\) for \(f_1,f_2,a,b,A,B\). Independently, a second checker
uses exact Gaussian-rational evaluation on a \(6\times6\) Cartesian grid.
The difference polynomial has degree at most five in each variable, so this
grid check is an exact interpolation proof, not merely numerical evidence;
four further complex points also pass. The two algebraic identities were
also rederived independently by hand.

For the alternative entire construction, eight selected complex points,
including both removable values, give maximum floating-point residuals
\(2.220\times10^{-16}\) for the row identity and
\(4.965\times10^{-16}\) for the Bezout identity. These are only implementation
checks. The proof of removability and the uniform growth estimate above,
independently rechecked, supplies the all-point justification. The scripts and
outputs are included with the chapter.

\paragraph{What remains uncertain.}
There is no remaining existence obstruction for the complex question just
formulated. The unanswered issue is documentary: the notes do not define the
index, clarify the coefficient field after the real example, or explain what
meaning of common divisor was intended. The complex polynomial certificate
invalidates the parenthesis under unrestricted complex coefficients; the real
degree argument invalidates its continuous-function half if all functions
remain real-valued. Neither fact licenses a reconstruction of the speaker's
missing restrictions. No assertion that a new open mathematical problem has
been solved, and no recommendation for a further research run, follows from
this critical verification.
